\documentclass[10pt]{article}
% Math Packages
\usepackage{amsmath, mathtools}
\usepackage{amssymb}
\usepackage{amsthm}
\usepackage{amsfonts}
\usepackage{bbm}
\usepackage{breqn}
\usepackage[margin=1in]{geometry}
\usepackage{graphicx}
\usepackage{tikz}
\usetikzlibrary{arrows.meta}
\usetikzlibrary{calc}
\usepackage{forest}
\usepackage{tikz-qtree}
\graphicspath{ {./images/} }
\usepackage{hyperref}
\usepackage[capitalize]{cleveref}
\usepackage[shortlabels]{enumitem}
\usetikzlibrary{arrows,matrix,positioning}
\usepackage{multicol}

\usepackage{listings}
\usepackage{xcolor}

\crefname{problem}{Problem}{Problems}
\Crefname{problem}{Problem}{Problems}

\lstset{
  language=R,
  basicstyle=\ttfamily\footnotesize,
  keywordstyle=\color{blue},
  commentstyle=\color{gray},
  stringstyle=\color{red!70!black},
  framesep=2pt,
  framexleftmargin=0pt, 
  frame=single,
  breaklines=false,
  showstringspaces=false
}
% for the pipe symbol
\usepackage[T1]{fontenc}

% Citing theorems by name. (source: https://tex.stackexchange.com/questions/109843/cleveref-and-named-theorems)
\makeatletter
\newcommand{\ncref}[1]{\cref{#1}\mynameref{#1}{\csname r@#1\endcsname}}

\def\mynameref#1#2{%
  \begingroup
  \edef\@mytxt{#2}%
  \edef\@mytst{\expandafter\@thirdoffive\@mytxt}%
  \ifx\@mytst\empty\else
  \space(\nameref{#1})\fi
  \endgroup
}
\makeatother

% Colorful Notes
\usepackage{color} \definecolor{Red}{rgb}{1,0,0} \definecolor{Blue}{rgb}{0,0,1}
\definecolor{Purple}{rgb}{.5,0,.5} \def\red{\color{Red}} \def\blue{\color{Blue}}
\def\gray{\color{gray}} \def\purple{\color{Purple}}
\newcommand{\rnote}[1]{{\red [#1]}} % \rnote{foo} gives '[foo]' in red
\newcommand{\pnote}[1]{{\purple [#1]}} % \pnote{foo} gives '[foo]' in purple
\newcommand{\bnote}[1]{{\blue #1}} % \bnote{foo} gives 'foo' in blue
\newcommand{\gnote}[1]{{\gray #1}} % \gnote{foo} gives 'foo' in gray
\newcommand{\Max}[1]{{\purple [#1]}} % \bnote{foo} then 'foo' is blue


% Claim numbering (the counter restarts after each proof environment)
\newcounter{claimcount}
\setcounter{claimcount}{0}
\newenvironment{claim}{\refstepcounter{claimcount}\par\addvspace{\medskipamount}\noindent\textbf{Claim \arabic{claimcount}:}}{}
\usepackage{etoolbox}
\AtBeginEnvironment{proof}{\setcounter{claimcount}{0}}
\newenvironment{claimproof}{\par\addvspace{\medskipamount}\noindent\textit{Proof of Claim  \arabic{claimcount}.}}{\hfill\ensuremath{\qedsymbol} \tiny{Claim}

  \medskip}
% Add claim support to cleverref
\crefname{claimcount}{Claim}{Claims}


% Math Environments
\newtheorem{theorem}{Theorem}
\newtheorem{assumption}[theorem]{Assumption}
\newtheorem{lemma}[theorem]{Lemma}
\newtheorem{proposition}[theorem]{Proposition}
\newtheorem{corollary}[theorem]{Corollary}
\newtheorem{question}[theorem]{Question}
\theoremstyle{definition}
\newtheorem*{definition}{Definition}
\newtheorem{remark}[theorem]{Remark}
\newtheorem{example}[theorem]{Example}
\newtheorem{notation}[theorem]{Notation}
\newtheorem{problem}[theorem]{Problem}

% Matrices and Column Vectors. 
\usepackage{stackengine}
\setstackgap{L}{1.0\normalbaselineskip}
\usepackage{tabstackengine}
\setstackEOL{;}% row separator
\setstackTAB{,}% column separator
\setstacktabbedgap{1ex}% inter-column gap
\setstackgap{L}{1.5\normalbaselineskip}% inter-row baselineskip
\let\nmatrix\bracketMatrixstack  %Usage: \nmatrix{1,2,3\4,5,6}
\newcommand\cv[1]{\setstackEOL{,}\bracketMatrixstack{#1}} %usage: \cv{1,2,3}

% Custom Math Coqmmands
\newcommand{\vt}{\vskip 5mm} % vertical space
\newcommand{\fl}{\noindent\textbf} % first line
\newcommand{\Fl}{\vt\noindent\textbf} % first line with space above
\newcommand{\norm}[1]{\left\lVert#1\right\rVert} % norm
\newcommand{\pnorm}[1]{\left\lVert#1\right\rVert_p} % p-norm
\newcommand{\qnorm}[1]{\left\lVert#1\right\rVert_q} % q-norm
\newcommand{\1}[1]{\textbf{1}_{\left[#1\right]}} % indicator function 
\def\limn{\lim_{n\to\infty}} % shortcut for lim as n-> infinity
\def\sumn{\sum_{n=1}^{\infty}} % shortcut for sum from n=1 to infinity
\def\sumkn{\sum_{k=1}^{n}} % shortcut for sum from k=1 to n
\def\sumin{\sum_{i=1}^{n}} % shortcut for sum from i=1 to n
\def\SAs{\sigma\text{-algebras}} % shortcut for $\sigma$-algebras
\def\SA{\sigma\text{-algebra}} % shortcut for $\sigma$-algebra
\def\Ft{\mathcal{F}_t} % time-indexed sigma-algebra (t)
\def\Fs{\mathcal{F}_s} % time-indexed sigma-algebra (s)
\def\F{\mathcal{F}} % sigma-algebra
\def\G{\mathcal{G}} % sigma-algebra
\def\R{\mathbb{R}} % Real numbers
\def\N{\mathbb{N}} % Natural numbers
\def\Z{\mathbb{Z}} % Integers
\def\E{\mathbb{E}} % Expectation
\def\P{\mathbb{P}} % Probability
\def\Q{\mathbb{Q}} % Q probability
\def\dist{\text{dist}} %Text 'dist' for things like 'dist(x,y)'
\newcommand{\indep}{\perp \!\!\! \perp}  %independence symbol
\def\Var{\mathrm{Var}} % Variance
\def\tr{\mathrm{tr}} % trace

% Brackets and Parentheses
\def\[{\left [}
    \def\]{\right ]}
% \def\({\left (}
%   \def\){\right )}



\usepackage{color}
\definecolor{Red}{rgb}{1,0,0}
\definecolor{Blue}{rgb}{0,0,1}
\definecolor{Purple}{rgb}{.5,0,.5}
\def\red{\color{Red}}
\def\blue{\color{Blue}}
\def\gray{\color{gray}}
\def\purple{\color{Purple}}
\definecolor{RoyalBlue}{cmyk}{1, 0.50, 0, 0}
\newcommand{\dempfcolor}[1]{{\color{RoyalBlue}#1}} 
\newcommand{\demph}[1]{\textcolor{RoyalBlue}{\textbf{\slshape #1}}} % Slanted

\usepackage{emoji}
% comment exactly one of the following line to show / hide the solutions
% \newcommand{\solution}[1]{{\purple #1}} % uncomment to show the solutions
\newcommand{\solution}[1]{} % uncomment to hide the solutions




\begin{document}
\begin{center}
  \section*{Math 472: Homework 07}
  \textit{Due Monday, March 23}
\end{center}


\begin{problem}[Exercise 9.8]
  Let  $Y_{1},\ldots,Y_{n}$ denote a random sample from a probability density
  function $f(y)$, which has unknown parameter $\theta$. If $\widehat{\theta}$
  is an unbiased estimator of $\theta$, then under very general conditions
  \begin{equation*}
    \Var(\widehat{\theta}) \geq \frac{1}{\mathcal{I}(\theta)},
    \quad \text{where} \quad 
    \mathcal{I}(\theta)= n \E\left[- \frac{\partial^{2}}{\partial \theta^{2} }\Big[\log(f(Y)) \Big]  \right].
  \end{equation*}
  (This is known as the Cramer-Rao inequality.) If
  $\Var(\widehat{\theta})=\frac{1}{\mathcal{I}(\theta)}$, the estimator
  $\widehat{\theta}$ is said to be \textit{efficient}.

  Suppose that $p(y)$ is the Poisson probability function with mean
  $\lambda>0$. Show that $\overline{Y}$ is an efficient estimator of $\lambda$. 
\end{problem}




\begin{problem}
  \label{problem:increasing-functions-and-maximizers}
  Given a function $h:\R \to \R$, we say that $x^{*}\in \R$ is a \demph{maximizer}
  of $h$ if
  \begin{equation*}
    h(x) \leq h(x^{*}) \text{ for all }x\in \R.
  \end{equation*}
  We say that $x^{*}\in \R$ is a \demph{local maximum} of $h$ if there exists
  some $\delta>0$ such that
  \begin{equation*}
    h(x) \leq h(x^{*}) \text{ for all }x\in (x^{*}-\delta,x^{*}+\delta).
  \end{equation*}
  We say that a function $h$ is \demph{strictly increasing} if
  \begin{equation*}
    h(x) < h(y) \text{ whenever } x<y.
  \end{equation*}
  Let $f:\R \to \R$ be a function, and let $g = \phi\circ f$, where $\phi$ is
  a strictly increasing function whose domain contains the range of $f$.
  \begin{enumerate}[(a)]
    \item \label{item:3} Show that $x^{*}$ is a maximizer of $f$ if and only
    if it is a maximizer of $g$. (You don't need to assume that $f$ or $\phi$
    is differentiable.)
    \item Show that $x^{*}$ is a local maxium of $f$ if and only if it is a
    local maximum of $g$. Conclude that $f$ and $g$ have the same local
    maxima. (Again, you don't need to assume that $f$ or $\phi$ is
    differentiable).
    \item Does the same conclusion hold for local minima? \textit{[Hint:
      consider the function $-f(x)$.]}
  \end{enumerate}
\end{problem}

In our class, when finding the maximum likelihood estimator, we need to
maximize the likelihood $L(\theta)$. The previous problem shows that it is
enough to maximize $\log L(\theta)$, since $\phi(x):=\log(x)$ is strictly
increasing. In fact, we can obtain stronger conclusions if we are allowed to
assume that $f$ and $\phi$ are differentiable, as shown in the next problem.

\begin{problem}
  \label{problem:more-increasing-function-stuff}
  Let $f,g,$ and $\phi$ be defined as in
  \cref{problem:increasing-functions-and-maximizers}. Further, assume that $f$
  and $\phi$ are both twice differentiable, and that $\phi'(x)>0$ for all
  $x\in \R$.
  \begin{enumerate}[(a)]
    \item \label{item:1} Show that $x^{*}\in \R$ is a critical point of $f$ if
    and only if it is a critical point of $g$.
    \item \label{item:2} Suppose that $x^{*}\in \R$ is a critical point of $f$
    and $g$. Show that $f''(x^{*})$ and $g''(x^{*})$ are either (i) both
    positive, (ii) both negative, or (iii) both zero.
    \item What additional information does part \ref{item:2} tell you about
    functions $f$ and $g$?
  \end{enumerate}
\end{problem}

\begin{problem}
  There is a subtlety to the assumptions used in
  \cref{problem:increasing-functions-and-maximizers,problem:more-increasing-function-stuff},
  which we highlight in this problem. Let $\phi:\R \to \R$ be a differentiable
  function.
  \begin{enumerate}[(a)]
    \item \label{item:4} Show that if $\phi'(x)>0$ for all $x\in \R$, then
    $\phi$ is strictly increasing. \textit{[Hint: use the mean value
      theorem]}.
    \item Show that the converse to part \ref{item:4} does not hold (i.e.,
    give a counterexample).
  \end{enumerate}
\end{problem}

\newpage
\begin{problem}
  Write \texttt{R} code to replicate Figure 9.2 (page 449) in the textbook.
\end{problem}

\begin{problem}[Exercise 9.3 and 9.15?]
  Let $Y_{1},\ldots,Y_{n}$ denote a random sample fro mthe uniform
  distribution on the interval $(\theta,\theta+1)$. Let
  \begin{equation*}
    \widehat{\theta}_{1} 
    = \overline{Y}-\frac{1}{2} \quad \text{and} \quad \widehat{\theta}_{2} 
    = \max(Y_{1},\ldots,Y_{n}) - \frac{n}{n-1}.
  \end{equation*}
  \begin{enumerate}[(a)]
    \item Show that both $\widehat{\theta}_{1}$ and $\widehat{\theta}_{2}$ are
    unbiased estiamtors of $\theta$.  
    \item Find the efficiency of $\hat{\theta}_{1}$ relative to
    $\widehat{\theta}_{2}$.
    \item Show that both $\hat{\theta}_{1}$ and $\widehat{\theta}_{2}$ are
    consistent estimators of $\theta$.   
  \end{enumerate}
\end{problem}

\begin{problem}
  Let $X$ be a random variable with cdf
  \begin{equation*}
    F(x) = \frac{1}{2}+ \frac{1}{\pi}\arctan(x), \quad (x\in \R)
  \end{equation*}
  (This is called a Cauchy random variable.)
  \begin{enumerate}[(a)]
    \item Show that $\E\left[X \right]$ is not defined.
    \item Let $Y_{n} = \frac{X_{1}+\ldots+X_{n}}{n}$. For
    $n\in \left\{100, 10000, 100000\right\}$, sample a large number of
    independent replicates of $Y_{n}$, and plot the histograms. Explain what
    you see.
  \end{enumerate}
\end{problem}


\begin{problem}[Exercise 9.37]
  Let $X_{1},\ldots,X_{n}$ denote $n$ independent and identically distributed
  Bernoulli random variables such that
  \begin{equation*}
    \P\left[X_{i}=1 \right] = p \quad \text{and} \quad  \P\left[ X_{i}=0\right]=1-p
  \end{equation*}
  for each $i=1,\ldots,n$. Let $S_{n} = X_{1}+\ldots+X_{n}$. Show that $S_{n}$
  is sufficient for $p$ by using the factorization criterion.
\end{problem}

\begin{problem}
  Suppose $Y_{1},\ldots,Y_{n}$ is a random sample from a geometric
  distribution with parameter $p$. Show that $\overline{Y}$ is sufficeient for
  $p$.
\end{problem}

\begin{problem}
  Let $Y_{1},\ldots,Y_{n}$ be a random sample from a normal distribution with
  mean $\mu$ and variance $\sigma^{2}$.
  \begin{enumerate}[(a)]
    \item If $\mu$ is unknown and $\sigma^{2}$ is known, show that
    $\overline{Y}$ is sufficient for $\mu$. 
    \item If $\mu$ is known and $\sigma^{2}$ is unknown, show that
    $\sum_{i=1}^{n}(Y_{i}-\mu)^{2}$ is sufficient for $\sigma^{2}$.
  \end{enumerate}
\end{problem}


\begin{problem}
  Suppose $Y_{1},\ldots,Y_{n}$ is a random sample from a Poisson distribution
  with mean $\lambda>0$.
  \begin{enumerate}[(a)]
    \item Use the factorization criterion to find a sufficient statistic for
    $\lambda$.
    \item Find the MLE $\widehat{\lambda}$ of $\theta$.
  \end{enumerate}
\end{problem}

\begin{problem}
  Suppose $Y_{1},\ldots,Y_{n}$ is a random sample from a uniform distribution
  with pdf
  \begin{equation*}
    f(y\mid \theta) = \frac{1}{2\theta+1}\cdot \mathbf{1}_{\left[ 0 \leq y \leq 2\theta+1\right]}.
  \end{equation*}
  Find the mle of $\theta$. 
\end{problem}

\end{document}